\section{Билет 18. Математическое описание волны. Уравнение плоской волны. График плоской волны. Параметры волны (длина волны, волновое число). Скорость распространения гармонических волн, ее связь с параметрами волны.}

\begin{definition}
Для описания волнового движения необходимо найти смещение $\xi$ частиц среды как функцию координаты $x$ и времени $t$. Рассмотрим простейший случай --- \textbf{плоскую монохроматическую волну}. Пусть источник колебаний находится в начале координат ($x=0$) и колеблется по закону:
\begin{equation}
    \xi(0, t) = A \cos(\omega t).
\end{equation}
Волна распространяется вдоль оси $x$ со скоростью $v$. Колебания дойдут до точки с координатой $x$ с временной задержкой $\tau = x/v$. Поскольку среда однородна и поглощением энергии пренебрегаем, амплитуда $A$ и частота $\omega$ сохраняются. Следовательно, смещение частиц в точке $x$ в момент $t$ равно смещению источника в более ранний момент $t - \tau$:
\begin{equation}
    \xi(x, t) = A \cos\left[\omega \left(t - \frac{x}{v}\right)\right]. \label{eq:plane_wave}
\end{equation}
Это \textbf{уравнение плоской бегущей волны}, распространяющейся в направлении возрастания $x$. Для волны, бегущей в противоположном направлении, знак меняется на плюс: $\xi(x, t) = A \cos[\omega(t + x/v)]$.
\end{definition}

\begin{property}[Параметры волны и скорость распространения]
\begin{itemize}
    \item \textbf{Период волны $T$} --- время одного полного колебания частиц среды.
    \item \textbf{Частота волны $\nu$} --- число колебаний в единицу времени: $\nu = 1/T$.
    \item \textbf{Длина волны $\lambda$} --- расстояние между ближайшими точками, колеблющимися в одинаковой фазе. За период $T$ волна распространяется на расстояние $\lambda$.
\end{itemize}
\textbf{Фазовая скорость} распространения волны связана с этими параметрами соотношениями:
\begin{equation}
    v = \frac{\lambda}{T} = \lambda \nu. \label{eq:wave_speed}
\end{equation}
Учитывая связь циклической частоты и периода $\omega = 2\pi/T$, уравнение волны \eqref{eq:plane_wave} можно переписать в виде:
\begin{equation}
    \xi(x, t) = A \cos\left(2\pi \frac{t}{T} - 2\pi \frac{x}{\lambda}\right) = A \cos(\omega t - kx). \label{eq:wave_k}
\end{equation}
\end{property}

\begin{definition}
Величина $k$, входящая в уравнение \eqref{eq:wave_k}, называется \textbf{волновым числом}:
\begin{equation}
    k = \frac{2\pi}{\lambda} = \frac{\omega}{v}.
\end{equation}
Волновое число показывает, сколько длин волн укладывается на отрезке длиной $2\pi$ метров. Единица измерения --- радиан на метр (рад/м).
\end{definition}

\begin{remark}
\textbf{График плоской волны.} Зависимость $\xi(x, t_0)$ при фиксированном моменте времени $t_0$ представляет собой синусоиду, показывающую распределение смещений \textbf{всех} частиц среды вдоль направления распространения волны в данный момент. Важно отличать его от графика колебаний $\xi(x_0, t)$, который показывает изменение смещения \textbf{одной} конкретной частицы во времени. Фазовая скорость $v = dx/dt$ характеризует скорость перемещения конкретной фазы колебаний (например, гребня волны) в пространстве.
\end{remark}
